\documentclass{article}

\usepackage{geometry}
\usepackage[T1]{fontenc}
\usepackage[OT1]{fontenc}
\usepackage{mathtools}
\usepackage{amsmath}
\usepackage{amsfonts}
\usepackage{amssymb}
\usepackage{textcomp}

\usepackage{calc}
\usepackage{caption}

\usepackage{litcode}
\usepackage{verbatim}

\def\codebg{yellow!10}
\def\vcodebg{green!10}

%%%%%%%%%%%%%%%%%%%%%%%%%%%%%%%%%%%%%%%%%%%%%%%%%%%%%%
\begin{document}

\title{Document 2}
\author{Praneet Kapor}
\date{\today}
\maketitle

\section{Sample 1}
I am going to present two examples to the reader. This is the first one, which will print \tt Hello, world!
\rm 100 times.

\begin{code}[main.c]
@<$\langle$\textit{main.c}$\rangle \equiv$@>
#include <stdio.h>
#include <stdlib.h>

int main(int argc, char **argv)
{
    @<\coderef{Print `Hello World' 100 times}@>
    return 0;
}
\end{code}

The following piece of code will print \tt Hello, world! \rm 100 times.

\begin{code}[Print `Hello World' 100 times]
@<$\langle$\textit{Print `Hello World' 100 times}$\rangle \equiv$@>
for (int i = 0; i < 100; ++i)
{
    fprintf(stdout, "Hello, world!\n");
}
\end{code}

Now here is some regular text.
\begin{code}
A big brown fox jumped over a little lazy dog.
You escape in code by using @<@<@> and @<@>@> as escape characters.
\end{code}

A more appropriate choice sometimes would be the \tt vcode \rm environment. It won't get captioned or
labelled.
\begin{vcode}
A big brown fox jumped over a little lazy dog.
You escape in vcode by using @< and @> as escape characters.
\end{vcode}

\lstlistoflistings
\end{document}

%%%%%%%%%%%%%%%%%%%%%%%%%%%%%%%%%%%%%%%%%%%%%%%%%%%%%%%%%